% =======================================================
% 8. DISCUSSION AND THREATS TO VALIDITY
% =======================================================
\section{Discussion and Threats to Validity}
\label{sec:discussion}

\paragraph{What the model contributes, and what it must not decide} The data support a division of labour rather than a ranking. On loops whose independence is a matter of \emph{recognizing intent}---a stencil written with flattened indices, a gather through a lookup table, a temporary declared outside the loop, a DOALL loop over pointer parameters that are never overlapped---current open-weight models are far more complete than a sound compiler analysis. On loops whose legality depends on \emph{facts outside the text of the loop}---which arguments a caller passes, whether an index array is a permutation---they answer from the loop body alone and are confidently wrong. The static veto handles the second category when the fact is a proven dependence, and the call-site guard handles it when the fact is visible at a call in the same translation unit. Facts that are invisible to both (an index array that is a permutation only by construction, a pointer passed across translation units) remain the residual risk of any hybrid scheme, and are the reason we keep the threshold and the veto explicit and auditable.

\paragraph{Confidence is a weak gate} Self-reported confidence separated correct from incorrect answers only partially (Table~\ref{tab:aidist}), and for GPT-OSS-120B the compiler facts lowered the reported confidence without improving correctness. A threshold must therefore be calibrated on held-out data for each model and prompt; an evaluation that tunes it on the test set, as our preliminary version did, overstates both safety and recall.

\paragraph{Practical deployment} In a CI/CD setting HF-Static runs at compiler speed, so the cost of the pipeline is dominated by the model calls on \textsc{Unknown} loops; on realistic C++ these are about half of all loops (Table~\ref{tab:pipeline}). The measured cost is a fraction of a cent per loop, and exact-match caching of answers for unchanged code is safe, but similarity-based caching should not be used for legality decisions. Since the harness exposed both a miscompilation in the toolchain and a soundness bug in our own analysis, every directive proposed by any tool should be validated by differential execution before it is merged; for kernel-sized code this is inexpensive.

\paragraph{Limitations of HF-Static} The analysis is intra-procedural (apart from the call-site guard), works on optimized IR rather than on source, and inherits the precision of LLVM's DependenceAnalysis, which does not delinearize flattened subscripts. Delinearization, polyhedral dependence testing on static control parts and inter-procedural mod/ref information would move loops from \textsc{Unknown} to \textsc{Parallel} and reduce the work left to the model. The call-site guard vetoes legal calls with disjoint halves of one buffer; a range-based overlap test would remove these false vetoes.

\subsection{Threats to validity}
\label{sec:threats}

\paragraph{Construct validity} Our legality definition (one loop-level directive, unchanged body) is stricter than ``can be parallelized after restructuring''. Some loops labelled serial admit a parallel scan, a loop-nest transformation or run-time alias checks; a tool that performs such rewrites (Polly transforms six of these nests) is penalized by our metrics. Conversely, a parallel label is supported by 12 executions, not by a proof; the schedules were chosen to maximize interleavings, and every canonical serial loop of HF-LoopBench did produce a deviating result under naive parallelization, which indicates that the harness is sensitive to the dependences it is meant to expose. The arg-min loops of HF-LoopBench-X are serial by definition but their race did not manifest.

\paragraph{Internal validity} Ground-truth labels were written by the authors, and the benchmarks were designed together with the tools. To limit this bias, the HF-Static rules and prompts were fixed before the LLM experiments and not changed after observing model outputs, labels were validated by execution, the threshold was selected on a development split disjoint (by base kernel) from the reported test split, and HF-LoopBench-X was written after all of this was frozen. The only change after freezing is the soundness fix of Section~\ref{sec:rq-ext}, which can only turn verdicts into \textsc{Serial}; we report its effect on every corpus. LLM outputs are not fully deterministic even at temperature~0; we quantify run-to-run agreement and release the complete request/response log, from which every number is recomputed.

\paragraph{External validity} The benchmarks consist of small kernels with one target loop each and generated variants; real applications contain larger functions, macros, C++ abstractions and inter-procedural effects. The 200-file corpus shows that the pipeline scales to such code but has no ground truth. The LLM study is limited by free-tier quotas and to three open-weight models available on one provider at the time of the study; the Llama~3.3~70B model of our preliminary version is no longer served there. All timings come from one laptop-class processor.
